\scanpage{84}
{\itshape
% letterhead
UNIVERSITE DES SCIENCES\\
ET TECHNIQUES DU LANGUEDOC

M A T H E M A T I Q U E S

D.E.U.G. A Première Année

Option Mathématique

1976

\uline{Titre}\qquad Polyèdres réguliers

Ce ``cours'' est destiné aux étudiants disposés à trouver de l'amusement en faisant des choses mathématiques très élémentaires et ``visuelles'', pouvant servir d'introduction à plusieurs courants d'idées et types d'intuition importants de géométrie et de topologie $n$-dimensionnelle. Il ne sera d'aucune utilité pour la préparation d'autres certificats de mathématiques. Il n'y aura pas de travaux dirigés proprement dits, et les séances seront en partie consacrées à la réflexion collective sur des questions se présentant dans l'exposé, avec des volontaires venant éventuellement expliquer une démonstration au tableau. Pas d'exercices notés. Le cours est divisé en deux parties, qui seront développées simultanément, l'une par A. GROTHENDIECK, l'autre par C. CONTOU CARRERE. Ces parties sont dans une large mesure indépendantes l'une de l'autre. Il y sera question des choses suivantes :

I) \uline{Polyèdres} : Polygones plans, division des plans en deux parties (``th. de Jordan''), cas convexe. Polyèdres convexes, facettes du polyèdre convexe, combinatoire d'un polyèdre convexe. Dualité des polyèdres convexes ayant un point intérieur donné. Peut-être des notions sur les polyèdres pas nécessairement convexes.

II) \uline{Les cinq polyèdres réguliers} (``solides platoniciens''). Description métrique, description combinatoire, éventuellement description ``projective'' en termes de ``configurations'' de points, droites et plans. Etude de leurs groupes de symétrie, liens avec les ``groupes symétriques'' $\handwritten{\unsure{\mathfrak S|\mathscr S}}_3$, $\handwritten{\mathfrak S}_4$, $\handwritten{\mathfrak S}_5$ des permutations d'ensembles finis : 3,4,5 objets, et leurs sous-groupes.

On présuppose une familiarité raisonnable de l'étudiant avec la géométrie du plan et de l'espace, qu'on abordera dans des points de vue nouveaux pour lui, et avec la notion de groupe, dont le ``cours'' permettra une compréhension plus approfondie par l'examen détaillé d'exemples concrets.
}
% archive sheet number: 83
